\begin{notechapter}
\chapter{Gradient Boosting}\label{ch:gradient-boosting}
\chaptersubtitle{Gradient descent in a space of functions}

\begin{chapterintro}
AdaBoost minimizes the exponential loss one term at a time
(\cref{ensemble-methods:sec:adaboost}). Gradient boosting (Friedman, \emph{Greedy Function
Approximation: A Gradient Boosting Machine}, Annals of Statistics, 2001) does the same
for any differentiable loss, by viewing the prediction function itself as the variable
of a gradient descent. Each step computes the negative gradient of the loss at the
training points --- the \emph{pseudo-residuals} --- and fits a tree to them. This
chapter makes the functional gradient precise, derives the algorithm and its
second-order variant used by XGBoost and LightGBM, and states what is known about its
convergence. You will learn
\begin{itemize}[nosep]
  \item the functional derivative of the empirical risk, and why its gradient depends on the
        inner product chosen on functions;
  \item why the raw gradient step would only memorize the training set, and how fitting a
        tree to the pseudo-residuals extends it;
  \item the algorithm with line search and shrinkage, and the pseudo-residuals of the
        common losses;
  \item Newton boosting, its optimal leaf values $-G/(H + \lambda)$ and split gain;
  \item what convergence and consistency results say, and the role of early stopping.
\end{itemize}
\prerequisites{\cref{ch:gradient-tangent-space} (the gradient depends on the metric);
\cref{ch:decision-trees,ch:ensemble-methods}; \cref{ch:output-distributions} for the
losses.}
\end{chapterintro}

\section{The problem}\label{gradient-boosting:sec:problem}

Given data $(x_i, y_i)_{i=1}^n$ and a loss $L(y, F)$ differentiable in its second argument, we look
for a function $F$ minimizing the empirical risk
\[
  \mathcal R[F] = \sum_{i=1}^nL\bigl(y_i, F(x_i)\bigr),
\]
an approximation of the population risk $\E[L(Y, F(X))]$. Here the variable is the function $F$
itself, not a parameter vector. $\mathcal R$ is a \emph{functional}: it maps functions to numbers.

\section{The functional gradient}\label{gradient-boosting:sec:functional}

\paragraph{Directional derivative.} The Gâteaux derivative of $\mathcal R$ at $F$ in the direction of a
function $h$ is
\[
  \mathcal R'[F; h] = \lim_{\epsilon\to0}\frac{\mathcal R[F + \epsilon h] - \mathcal R[F]}\epsilon = \sum_{i=1}^ng_i\,h(x_i),\qquad
  g_i = \frac{\partial L(y_i, F)}{\partial F}\Big|_{F = F(x_i)},
\]
by the chain rule applied to each term. It is linear in $h$, and it depends on $h$ only through
its values at the training points: $\mathcal R'[F;\cdot] = \sum_ig_i\,\delta_{x_i}$, a combination of the point
evaluations $\delta_{x_i}: h\mapsto h(x_i)$ (the ``Dirac deltas'' at the training points).

\paragraph{From derivative to gradient.} As in \cref{ch:gradient-tangent-space}, a derivative
is a linear form; a \emph{gradient} is the vector that represents it through an inner
product, $\mathcal R'[F; h] = \inner{\nabla\mathcal R[F]}h$, and it depends on the inner product chosen. Three
choices on spaces of functions are instructive.
\begin{itemize}
  \item $L^2(\mu)$, with $\inner fg = \int fg\,d\mu$ for the input distribution $\mu$. If $\mu$ has a density, no function
        represents $h\mapsto\sum_ig_ih(x_i)$: changing $h$ at the finitely many points $x_i$ changes the
        derivative but not any integral. The linear form is not continuous on $L^2(\mu)$ and there
        is no gradient.
  \item $L^2(\hat\mu_n)$, the functions on the training points with $\inner fg_n = \frac1n\sum_if(x_i)g(x_i)$. Here
        the gradient exists and is the vector of values $(ng_1,\dots,ng_n)$; with the inner product
        $\sum_if(x_i)g(x_i)$ it is simply $(g_1,\dots,g_n)$.
  \item A reproducing kernel Hilbert space with kernel $k$ (\cref{ch:rkhs}), where point evaluation is
        continuous and represented by $k(x_i,\cdot)$: $h(x_i) = \inner h{k(x_i,\cdot)}$. The gradient is then the smooth
        function $\sum_ig_i\,k(x_i,\cdot)$, defined everywhere. Gradient descent in this space is
        \emph{kernel boosting}, and every iterate is a kernel expansion, as the representer
        theorem predicts.
\end{itemize}
With the second choice, the steepest-descent step changes $F$ only at the training points:
\[
  F_{m+1}(x_i) = F_m(x_i) - \rho\,g_i,\qquad F_{m+1}(x) = F_m(x)\ \text{elsewhere}.
\]
Iterated, it drives the training loss to its infimum by memorizing the training set, and
it says nothing about new inputs.

\paragraph{Restricting the direction.} Gradient boosting chooses the descent direction
in a class $\Hs$ of simple functions, typically regression trees of a fixed size. The best
approximation of the negative gradient in $\Hs$, in the norm of $L^2(\hat\mu_n)$, is the least-squares
fit to the \emph{pseudo-residuals} $r_i = -g_i$:
\[
  h_m = \argmin_{h\in\Hs}\sum_{i=1}^n\bigl(r_i - h(x_i)\bigr)^2 .
\]
When $\Hs$ is a linear space this is the orthogonal projection of $-\nabla\mathcal R$ onto $\Hs$, and $h_m$ is the
steepest-descent direction within $\Hs$; for trees, which are not a linear space, it is the
greedy approximation computed by the tree-growing algorithm of
\cref{decision-trees:sec:greedy}. Either way, $h_m$ is defined on the whole input space: the tree
generalizes the gradient from the training points to new inputs.

\section{The algorithm}\label{gradient-boosting:sec:algorithm}

\begin{definition}[Gradient tree boosting]\label{gradient-boosting:def:algorithm}
\begin{enumerate}[nosep]
  \item Initialize with the best constant: $F_0 = \argmin_c\sum_iL(y_i, c)$.
  \item For $m = 1,\dots,M$:
  \begin{enumerate}[nosep]
    \item compute the pseudo-residuals $r_{im} = -\partial L(y_i, F)/\partial F$ at $F = F_{m-1}(x_i)$;
    \item fit a regression tree with leaves $R_{1m},\dots,R_{Jm}$ to the pairs $(x_i, r_{im})$ by least squares;
    \item for each leaf, find the best constant step,
          $\gamma_{jm} = \argmin_\gamma\sum_{x_i\in R_{jm}}L\bigl(y_i, F_{m-1}(x_i) + \gamma\bigr)$;
    \item update $F_m(x) = F_{m-1}(x) + \nu\sum_j\gamma_{jm}\bbone[x\in R_{jm}]$.
  \end{enumerate}
\end{enumerate}
\end{definition}

Step (c) is a line search, done separately in each leaf: the tree supplies the cells, the
loss decides the values. For the squared loss it gives the leaf means of the residuals,
the values the tree already has. The \emph{shrinkage} $0<\nu\leq1$ (the learning rate) takes only a
fraction of each step; values of $0.01$ to $0.1$ with many more trees generalize markedly
better than $\nu = 1$, an empirical fact that the analysis of early stopping explains in
part. \emph{Stochastic} gradient boosting (Friedman, 2002) fits each tree to a random half
or so of the data, which both speeds it up and regularizes.

\paragraph{Pseudo-residuals of the usual losses.}
\begin{center}
\small
\begin{tabular}{@{}lll@{}}
\toprule
loss & $L(y, F)$ & pseudo-residual $-\partial L/\partial F$ \\
\midrule
squared & $\frac12(y - F)^2$ & $y - F$ \\
absolute & $\abs{y - F}$ & $\sign(y - F)$ \\
Huber ($\delta$) & $\frac12(y - F)^2$ if $\abs{y - F}\leq\delta$, else $\delta(\abs{y - F} - \frac\delta2)$ & $y - F$, clipped to $[-\delta, \delta]$ \\
logistic, $y\in\{0, 1\}$ & $\log(1 + e^F) - yF$ & $y - \sigma(F)$ \\
exponential, $y\in\{\pm1\}$ & $e^{-yF}$ & $ye^{-yF}$ \\
\bottomrule
\end{tabular}
\end{center}
For the squared loss the pseudo-residuals are the residuals, and gradient boosting
repeatedly fits the errors of the current model. For the logistic loss, $F$ is the
log-odds and the pseudo-residual is the label minus the predicted probability
(\cref{sigmoid-softmax:sec:sigmoid}). For the absolute loss only the signs of the residuals
enter the tree, and the leaf values of step (c) are medians: the algorithm is robust to
outliers in $y$.

\begin{example}[Three rounds on eight points]\label{gradient-boosting:ex:rounds}
Data $x = 1,\dots,8$, $y = (1, 2, 2.5, 4, 3.5, 3, 1.5, 1)$, squared loss, stumps, $\nu = 0.5$. Then $F_0 = \bar y = 2.3125$,
with total squared error $8.969$.
\begin{center}
\begin{tabular}{@{}clccc@{}}
\toprule
$m$ & residuals $r_{im}$ & split & leaf values & SSE of $F_m$ \\
\midrule
1 & $(-1.31, -0.31, 0.19, 1.69, 1.19, 0.69, -0.81, -1.31)$ & $x\leq6.5$ & $0.354, -1.062$ & $6.711$ \\
2 & $(-1.49, -0.49, 0.01, 1.51, 1.01, 0.51, -0.28, -0.78)$ & $x\leq2.5$ & $-0.990, 0.330$ & $4.752$ \\
3 & $(-1.00, 0.01, -0.16, 1.35, 0.85, 0.35, -0.45, -0.95)$ & $x\leq6.5$ & $0.232, -0.696$ & $3.783$ \\
\bottomrule
\end{tabular}
\end{center}
The hump-shaped target is approached by alternately cutting its right and left
shoulders; the same split can be chosen again, since the residuals change. With $\nu = 0.5$
each tree contributes half its fit.
\end{example}

\section{Newton boosting}\label{gradient-boosting:sec:newton}

XGBoost (Chen and Guestrin, 2016) and LightGBM use a second-order expansion of the loss
around the current predictions. With $g_i$ and $h_i = \partial^2L/\partial F^2$ at $F_{m-1}(x_i)$, a new tree $f$ with leaf
values $w_1,\dots,w_J$ is chosen to minimize
\[
  \sum_i\Bigl(g_if(x_i) + \tfrac12h_if(x_i)^2\Bigr) + \gamma J + \tfrac\lambda2\sum_jw_j^2,
\]
where $\gamma$ penalizes the number of leaves and $\lambda$ the leaf values. Grouping the points by leaf,
with $G_j = \sum_{x_i\in R_j}g_i$ and $H_j = \sum_{x_i\in R_j}h_i$, the objective is $\sum_j(G_jw_j + \frac12(H_j + \lambda)w_j^2) + \gamma J$, a sum of
independent parabolas. Hence
\[
  w_j^* = -\frac{G_j}{H_j + \lambda},\qquad\text{objective}^* = -\frac12\sum_j\frac{G_j^2}{H_j + \lambda} + \gamma J,
\]
and a split of a leaf into $L$ and $R$ is worth
\[
  \mathrm{Gain} = \frac12\Bigl(\frac{G_L^2}{H_L + \lambda} + \frac{G_R^2}{H_R + \lambda} - \frac{(G_L + G_R)^2}{H_L + H_R + \lambda}\Bigr) - \gamma,
\]
which replaces the impurity decrease when growing the tree. Each leaf value is a Newton
step $-G/H$, shrunk by $\lambda$. For the squared loss ($h_i = 1$) and $\lambda = 0$ it is the mean residual; for
the logistic loss it is $\sum r_i/\sum p_i(1 - p_i)$, the leaf value Friedman already derived for
step (c).

\section{Convergence and regularization}\label{gradient-boosting:sec:theory}

\paragraph{Optimization.} For a convex loss, smooth in $F$, and a class $\Hs$ rich enough, gradient
boosting with suitable step sizes drives the empirical risk to its infimum. With
\emph{exact} functional gradient steps on a convex smooth risk whose infimum is attained,
the classical analysis of gradient descent gives an excess of order $C/M$ after $M$
iterations. With weak learners that only approximate the negative gradient, rates of
this type require further assumptions (each $h_m$ sufficiently correlated with the negative
gradient), and the known guarantees are weaker; schematically, the excess empirical risk
still decreases as $O(1/M)$ under such assumptions.

\paragraph{Statistics.} Driving the empirical risk to its infimum is not the goal: run to
convergence, boosting overfits. Zhang and Yu (\emph{Boosting with Early Stopping:
Convergence and Consistency}, Annals of Statistics, 2005) prove that it is
\emph{consistent} when stopped early: for a loss of the form $\phi(yF)$, base learners of
finite VC dimension (for example trees with a bounded number of leaves), restricted step
sizes $h_k$ with $\sum_kh_k = \infty$ and $\sum_kh_k^2<\infty$, and a number of iterations
$k_n\to\infty$ whose total step size $\sum_{k\leq k_n}h_k$ grows slowly enough with $n$, the
expected population risk converges to the best one in the span of the base learners,
\[
  \E_{\mathcal{D}_n}\bigl[\E[L(Y, F_{k_n}(X))]\bigr]\;\longrightarrow\;\inf_{F\in\operatorname{span}}\E[L(Y, F(X))],
\]
where the outer expectation is over the training sample $\mathcal{D}_n$. When that span is dense, the infimum is the
minimal risk. The number of iterations $M$ is thus the main regularization
parameter, chosen by validation (early stopping), together with $\nu$, the tree size and
the subsampling rate. Informally, stopping after $M$ small steps restricts $F$ to functions
reachable by a short path from $F_0$, a form of implicit regularization, as for early
stopping in neural networks. In bias--variance terms, each iteration lowers the bias,
and shrinkage, subsampling and small trees keep the variance in check.

\begin{chaptersummary}
\begin{itemize}[nosep]
  \item The derivative of the empirical risk is $h\mapsto\sum_ig_ih(x_i)$, a combination of point
        evaluations. Its gradient depends on the inner product: none in $L^2(\mu)$, the vector
        $(g_i)$ on the training points, $\sum_ig_ik(x_i,\cdot)$ in an RKHS.
  \item A raw gradient step only changes $F$ at the training points. Gradient boosting fits a
        tree to the pseudo-residuals $-g_i$ by least squares, then sets leaf values by line
        search and shrinks the step by $\nu$.
  \item Pseudo-residuals: residuals (squared loss), signs (absolute), $y - p$ (logistic).
  \item Newton boosting uses $g_i$ and $h_i$: leaf values $-G/(H + \lambda)$ and a split gain in closed form.
  \item Boosting run to convergence overfits; with early stopping it is consistent.
\end{itemize}
\end{chaptersummary}

\section*{Exercises}
\addcontentsline{toc}{section}{Exercises}

\begin{exercise}\label{gradient-boosting:exr:1}
Show that for the squared loss, step (c) of \cref{gradient-boosting:def:algorithm} returns the
mean residual of each leaf, and that for the absolute loss it returns a median of the
residuals of each leaf.
\solutionlink{gradient-boosting:exr:1}
\end{exercise}

\begin{exercise}\label{gradient-boosting:exr:2}
For the logistic loss with $y\in\{0, 1\}$, compute $g_i$ and $h_i$, and show that the Newton leaf value
with $\lambda = 0$ is $\sum_{i\in R}(y_i - p_i)/\sum_{i\in R}p_i(1 - p_i)$. What is $F_0$?
\solutionlink{gradient-boosting:exr:2}
\end{exercise}

\begin{exercise}\label{gradient-boosting:exr:3}
Show that gradient boosting with the exponential loss, $\nu = 1$ and weak learners with values in
$\{-1, 1\}$ chosen to maximize $\sum_ir_ih(x_i)$, followed by a line search, is AdaBoost.
\solutionlink{gradient-boosting:exr:3}
\end{exercise}

\begin{exercise}\label{gradient-boosting:exr:4}
Check the first row of \cref{gradient-boosting:ex:rounds}: compute the total squared error of the
residuals for the thresholds $2.5$ and $6.5$, and the leaf values of the best one.
\solutionlink{gradient-boosting:exr:4}
\end{exercise}
\end{notechapter}
